\chapter{Исследовательская часть}
\section{Технические характеристики}

Технические характеристики устройства, на котором выполнялись замеры времени, представлены далее.

\begin{enumerate}
	\item Процессор	AMD Ryzen 5 5500u ЦПУ 2592 МГц, ядер: 8, логических процессоров: 12.
	\item Оперативная память: 16 ГБ.
	\item Операционная система: Fedora 39.
\end{enumerate}

При замерах времени ноутбук был включен в сеть электропитания. Запущены были только системные приложения.

\section{Демонстрация работы программы}

На рисунке \ref{img:workExample} представлена демонстрация работы разработанного приложения для алгоритмов поиска.

\includeimage
{workExample}
{f}
{H}
{0.75\textwidth}
{Пример работы программы}

\section{Временные характеристики}
Результаты исследования замеров по времени приведены в таблице \ref{table:timings}.
Для чтения таблиц вводятся следующие обозначения:
\begin{enumerate}[label=\arabic*)]
	\item N~---~количество строк и столбцов умножаемых матриц;
	\item StandartMul~---~реализация стандартного алгоритма умножения матриц;
	\item StandartWinograd~---~реализация стандартного алгоритма умножения матриц Винограда;
	\item OptWinograd~---~реализация оптимизированного алгоритма умножения матриц Винограда.
\end{enumerate}

Для каждого N проводилось 1000 замеров времени, после чего результаты усреднялись. Все результаты вычислений приведены в миллисекундах.

\newpage
\begin{table}[!ht]
	\centering
	\caption{Таблица временных характеристик различных реализаций алгоритмов умножения матриц}
	\begin{tabular}{|c|c|c|c|c|}
		\hline
		N   & StandartMul (мс)  & StandartWinograd (мс) & OptWinograd (мс) \\ \hline
		1   & $2.48 \cdot 10^{-6}$   & $2.921 \cdot 10^{-6}$  & $2.901 \cdot 10^{-6}$   \\ \hline
		2   & $3.201 \cdot 10^{-6}$  & $3.69 \cdot 10^{-6}$   & $3.62 \cdot 10^{-6}$    \\ \hline
		3   & $4.36 \cdot 10^{-6}$   & $4.881 \cdot 10^{-6}$  & $4.6 \cdot 10^{-6}$     \\ \hline
		4   & $6.041 \cdot 10^{-6}$  & $6.251 \cdot 10^{-6}$  & $6.101 \cdot 10^{-6}$   \\ \hline
		5   & $8.502 \cdot 10^{-6}$  & $8.521 \cdot 10^{-6}$  & $7.962 \cdot 10^{-6}$   \\ \hline
		6   & $1.3013 \cdot 10^{-5}$ & $1.1192 \cdot 10^{-5}$ & $1.0442 \cdot 10^{-5}$  \\ \hline
		7   & $1.8053 \cdot 10^{-5}$ & $1.5502 \cdot 10^{-5}$ & $1.3673 \cdot 10^{-5}$  \\ \hline
		8   & $2.4444 \cdot 10^{-5}$ & $1.9473 \cdot 10^{-5}$ & $1.8113 \cdot 10^{-5}$  \\ \hline
		9   & $3.2396 \cdot 10^{-5}$ & $2.6394 \cdot 10^{-5}$ & $2.3203 \cdot 10^{-5}$  \\ \hline
		10  & $4.2668 \cdot 10^{-5}$ & $3.4646 \cdot 10^{-5}$ & $2.8644 \cdot 10^{-5}$  \\ \hline
		11  & $7.1692 \cdot 10^{-5}$ & $5.817 \cdot 10^{-5}$  & $4.5698 \cdot 10^{-5}$  \\ \hline
		12  & $6.8082 \cdot 10^{-5}$ & $5.3019 \cdot 10^{-5}$ & $4.9418 \cdot 10^{-5}$  \\ \hline
		13  & $8.4184 \cdot 10^{-5}$ & $6.6441 \cdot 10^{-5}$ & $5.8209 \cdot 10^{-5}$  \\ \hline
		14  & 0.000102997 & $7.8843 \cdot 10^{-5}$ & $7.1072 \cdot 10^{-5}$ \\ \hline
		15  & 0.000124351 & $9.7426 \cdot 10^{-5}$ & $8.4874 \cdot 10^{-5}$ \\ \hline
		20  & 0.000311187 & 0.000206624 & 0.000184481 \\ \hline
		25  & 0.000527658 & 0.000423793 & 0.000352759 \\ \hline
		30  & 0.00071889 & 0.000531073 & 0.000448624 \\  \hline
		35  & 0.000507514 & 0.000374357 & 0.000312252 \\ \hline
		40  & 0.000731114 & 0.000520318 & 0.000436362 \\ \hline
		45  & 0.00103469 & 0.000712128 & 0.000634552 \\ \hline
		50  & 0.00140134 & 0.000975394 & 0.000852307 \\ \hline
	\end{tabular}
	\label{table:timings}
\end{table}

Ниже представлен график, иллюстрирующий представленные в таблице \ref{table:timings} данные.

\begin{figure}[H]
	\centering
	\includesvg[inkscapelatex=false, width=0.8\textwidth, height=0.8\textheight]{inc/img/plot.svg}
	\caption{График сравнения по времени алгоритмов умножения матриц}
	\label{plot:plot}
\end{figure}

\section*{Вывод}
Из графика~\ref{plot:plot} сделан вывод, что среди анализируемых алгоритмов при больших размерностях матриц (более 5) самым эффективным является оптимизированный алгоритм Винограда, а стандартный алгоритм Винограда работает быстрее стандартного алгоритма умножения матриц.


